\begin{examplebox}
	\textbf{\textcolor{CExample}{例1（基础）}}

	\textbf{\textcolor{CExample}{题目：}}
	将函数 $y=\sin x$ 的图像，通过先平移后伸缩变换，得到 $y=3\sin(2x+\frac{\pi}{4})$。请写出具体的变换步骤。

	\textbf{\textcolor{CExample}{解题步骤：}}
	\begin{description}[style=nextline, leftmargin=2.8em, labelsep=0.5em, itemsep=0.45em]
		\item[\textbf{第一步（平移相位）：}] 将 $y=\sin x$ 向左平移 $\frac{\pi}{4}$ 个单位，得到 $y=\sin(x+\frac{\pi}{4})$。
			\par\textbf{理由：}
			$\varphi=\frac{\pi}{4}>0$，依据“左加右减”，向左移动 $|\varphi|$。
		\item[\textbf{第二步（横缩周期）：}] 将横坐标变为原来的 $\frac{1}{2}$，曲线被横向压缩，得到 $y=\sin(2x+\frac{\pi}{4})$。
			\par\textbf{理由：}
			$\omega=2$，横坐标乘以 $\frac{1}{\omega}=\frac{1}{2}$，即 $x$ 换为 $2x$。
		\item[\textbf{第三步（纵缩振幅）：}] 将纵坐标变为原来的 $3$ 倍，得到 $y=3\sin(2x+\frac{\pi}{4})$。
			\par\textbf{理由：}
			$A=3>0$，直接乘以 $|A|$。
	\end{description}

	\textbf{\textcolor{CExample}{关键结论：}}
	\textbf{本例 $A>0$，所以最后纵坐标乘以 $A=|A|=3$；一般情形应乘以带符号的 $A$。}

	\medskip
	\textbf{\textcolor{CExample}{例2（稍变形）}}

	\textbf{\textcolor{CExample}{题目：}}
	将 $y=\sin x$ 变换到 $y=2\sin(3x-\frac{\pi}{2})$，若按先伸缩后平移的顺序，求平移量及具体步骤。

	\textbf{\textcolor{CExample}{解题步骤：}}
	\begin{description}[style=nextline, leftmargin=2.8em, labelsep=0.5em, itemsep=0.45em]
		\item[\textbf{第一步（横缩周期）：}] 将横坐标变为原来的 $\frac{1}{3}$，得到 $y=\sin(3x)$。
			\par\textbf{理由：}
			$\omega=3$，横缩比例 $\frac{1}{\omega}=\frac{1}{3}$。
		\item[\textbf{第二步（右移相位）：}] 计算平移量：
			\[
				\begin{aligned}
					t &= \left|\frac{\varphi}{\omega}\right| \\
					&= \frac{\pi}{6}.
			\end{aligned}
			\]
			（其中 $\varphi=-\pi/2$, $\omega=3$）
			因 $\varphi<0$，应向右平移 $\pi/6$，得到
			\[
				\begin{aligned}
					y &= \sin\!\big[3\big(x-\frac{\pi}{6}\big)\big] \\
					&= \sin(3x-\frac{\pi}{2}).
			\end{aligned}
			\]
			\par\textbf{理由：}
			先横缩后，平移量调整为 $|\varphi/\omega|$，方向仍由 $\varphi$ 符号决定（负右移）。
		\item[\textbf{第三步（纵缩振幅）：}] 纵坐标变为原来的 $2$ 倍，得到 $y=2\sin(3x-\frac{\pi}{2})$。
			\par\textbf{理由：}
			$A=2$，纵缩 $|A|$ 倍。
	\end{description}

	\textbf{\textcolor{CExample}{关键结论：}}
	\textbf{先伸缩后平移的平移量是 $|\varphi/\omega|$，不是 $|\varphi|$。}

	\medskip
	\textbf{\textcolor{CExample}{例3（综合应用）}}

	\textbf{\textcolor{CExample}{题目：}}
	已知 $y=\sin x$，按下列两种变换顺序分别操作，分别写出得到 $y=4\sin(2x-\frac{\pi}{3})$ 的步骤，并比较两种路径的平移量。

	\textbf{\textcolor{CExample}{解题步骤：}}
	\begin{description}[style=nextline, leftmargin=2.8em, labelsep=0.5em, itemsep=0.45em]
		\item[\textbf{第一步（先移后缩）：}] 先右移 $\frac{\pi}{3}$，得 $y=\sin(x-\frac{\pi}{3})$；再横缩 $\frac{1}{2}$，得 $y=\sin(2x-\frac{\pi}{3})$；最后纵缩 $4$，得 $y=4\sin(2x-\frac{\pi}{3})$。
			\par\textbf{理由：}
			$\varphi=-\pi/3<0$，故右移 $|\varphi|=\pi/3$；$\omega=2$，横缩 $1/2$；$A=4$，纵缩 $4$。
		\item[\textbf{第二步（先缩后移）：}] 先横缩 $\frac{1}{2}$，得 $y=\sin(2x)$；再右移 $\frac{\pi}{6}$，得 $y=\sin[2(x-\frac{\pi}{6})]=\sin(2x-\frac{\pi}{3})$；最后纵缩 $4$，得 $y=4\sin(2x-\frac{\pi}{3})$。
			\par\textbf{理由：}
			先横缩后，平移量 $|\varphi/\omega| = (\pi/3)/2 = \pi/6$，方向向右。
		\item[\textbf{第三步（平移量对比）：}] 先移后缩平移量 $|\varphi|=\pi/3$，先缩后移平移量 $|\varphi/\omega|=\pi/6$，且 $|\varphi/\omega| = |\varphi|/\omega$。
			\par\textbf{理由：}
			$\omega=2$ 导致先缩后移的平移量减小一半。
	\end{description}

	\textbf{\textcolor{CExample}{关键结论：}}
	\textbf{两种变换路径都可以得到相同图像，但平移量不同；先伸缩后平移的平移量为 $|\varphi|/\omega$。}

	\medskip
	\textbf{\textcolor{CExample}{例4（负系数与翻折）}}

	目标为 $y=-2\sin(2x+\frac{\pi}{4})$。先移后缩：$\sin x\to\sin(x+\frac{\pi}{4})\to\sin(2x+\frac{\pi}{4})\to-2\sin(2x+\frac{\pi}{4})$，依次左移 $\pi/4$、横坐标乘以 $1/2$、纵坐标乘以 $-2$。先缩后移：$\sin x\to\sin(2x)\to\sin[2(x+\frac{\pi}{8})]\to-2\sin(2x+\frac{\pi}{4})$，其中第二步左移 $\pi/8$。最后一步也可分成纵坐标乘以 $2$，再关于 $x$ 轴对称；振幅是 $2$。代入 $x=\pi/8$，两条路径均得 $-2$，若只乘 $|A|$ 则误得 $+2$。

	\medskip
	\textbf{\textcolor{CExample}{例5（相等情形与横向拉伸）}}

	当 $\omega=1$、$\varphi=\pi/4$ 时，两条路径均左移 $\pi/4$，距离相等；当 $\varphi=0$、$\omega=2$ 时，两条路径均无水平平移。再取 $A=2$、$\omega=1/2$、$\varphi=-\pi/4$：先右移 $\pi/4$ 再把横坐标变为 $2$ 倍，得到 $\sin(x/2-\pi/4)$；或先把横坐标变为 $2$ 倍，再右移 $\pi/2$，得到同一函数，最后纵坐标乘以 $2$。这里 $0<\omega<1$ 对应横向\textbf{拉伸}，负 $\varphi$ 对应\textbf{右移}。
\end{examplebox}
